\documentclass[11pt]{article}

\usepackage[a4paper,margin=1in]{geometry}
\usepackage{amsmath,amssymb,amsthm,mathtools}
\usepackage{enumitem}
\usepackage{hyperref}

\newtheorem{theorem}{Theorem}
\newtheorem{lemma}[theorem]{Lemma}
\newtheorem{proposition}[theorem]{Proposition}
\newtheorem{remark}[theorem]{Remark}

\DeclareMathOperator{\diag}{diag}
\DeclareMathOperator{\sgn}{sgn}
\DeclareMathOperator{\dist}{dist}
\DeclareMathOperator{\spec}{spec}
\DeclareMathOperator{\nul}{null}
\DeclareMathOperator{\rank}{rank}
\DeclareMathOperator{\tr}{tr}

\begin{document}

\title{The Median Eigenvalue Branch and the Smallest Singular Value}
\author{}
\date{}
\maketitle

\section*{Statement}

Fix an integer $n\geq 2$ and a parameter $0<a<1$.  Let
\[
A_n(a)=
\begin{pmatrix}
0&1&&&0\\
a&0&1&&\\
& a&0&\ddots&\\
&&\ddots&\ddots&1\\
0&&&a&0
\end{pmatrix},
\]
and let $R_n$ be the reversal matrix,
\[
R_n e_j=e_{n+1-j},\qquad j=1,\ldots,n.
\]
For $x\in\mathbb R$ put
\[
B_n(x)=xI_n-A_n(a),\qquad C_n(x)=R_nB_n(x).
\]
The matrix $C_n(x)$ is real symmetric.  Let
\[
\lambda_1^{(n)}(x)\leq \lambda_2^{(n)}(x)\leq \cdots\leq \lambda_n^{(n)}(x)
\]
be the ordered eigenvalues of $C_n(x)$, and define the median index
\[
q_n=\left\lfloor\frac n2\right\rfloor+1.
\]
The eigenvalues of $A_n(a)$ are
\[
\mu_j^{(n)}=2\sqrt a\cos\frac{j\pi}{n+1},\qquad j=1,\ldots,n,
\]
ordered as
\[
\mu_1^{(n)}>\mu_2^{(n)}>\cdots>\mu_n^{(n)}.
\]
The goal is to prove
\[
\boxed{
 s_{\min}\bigl(xI_n-A_n(a)\bigr)
 =
 \left|\lambda_{q_n}^{(n)}(x)\right|
 }
\]
for every
\[
x\in[\mu_n^{(n)},\mu_1^{(n)}].
\]

\section*{A finite Green-kernel input}

The proof uses the following standard finite-dimensional form of the
Gantmacher--Krein oscillation theorem for Green kernels.  It is stated here
in the exact form needed below.

\begin{theorem}[Finite Green-kernel simplicity theorem]\label{thm:GK}
Let $0<s<1$, let $n\geq 2$, and let
\[
\theta\in\left(\frac{r\pi}{n+1},\frac{(r+1)\pi}{n+1}\right),
\qquad r=1,\ldots,n-1.
\]
Define the real symmetric matrix $K_{n,s,\theta}=(K_{ij})_{i,j=1}^n$ by
\[
K_{ij}
=
\frac{
 s^{i+j-n-2}\sin(i\theta)\sin(j\theta)
}{
 \sin\theta\,\sin((n+1)\theta)
},
\qquad i+j\leq n+1,
\]
and
\[
K_{ij}
=
\frac{
 s^{i+j-n-2}\sin((n+1-i)\theta)\sin((n+1-j)\theta)
}{
 \sin\theta\,\sin((n+1)\theta)
},
\qquad i+j\geq n+1.
\]
Then the spectral radius of $K_{n,s,\theta}$ is attained by a unique simple
real eigenvalue.  More precisely,
\[
(-1)^r\rho(K_{n,s,\theta})
\]
is a simple eigenvalue of $K_{n,s,\theta}$, and every other eigenvalue $\nu$
satisfies
\[
|\nu|<\rho(K_{n,s,\theta}).
\]
\end{theorem}

\begin{remark}
Theorem~\ref{thm:GK} is the finite discrete Sturm--Gantmacher--Krein theorem:
the Green kernel in the $r$-th spectral gap is strictly sign-regular, and the
variation-diminishing property gives a unique dominant eigenvalue with exactly
the corresponding oscillation class.  This is the only external oscillation
input in the argument.
\end{remark}

\section*{Step 1: Symmetrization}

\begin{lemma}\label{lem:symmetric}
For every $x\in\mathbb R$, the matrix
\[
C_n(x)=R_n(xI_n-A_n(a))
\]
is real symmetric, and
\[
s_j(C_n(x))=s_j(xI_n-A_n(a)),\qquad j=1,\ldots,n.
\]
Consequently,
\[
s_{\min}(xI_n-A_n(a))
=
\min_{1\leq j\leq n}\left|\lambda_j^{(n)}(x)\right|.
\]
\end{lemma}

\begin{proof}
Since
\[
A_n(a)^T=R_nA_n(a)R_n,
\]
we have
\[
C_n(x)^T
=
(xI_n-A_n(a))^TR_n
=(xI_n-A_n(a)^T)R_n
=R_n(xI_n-A_n(a))
=C_n(x).
\]
Thus $C_n(x)$ is real symmetric.  Since $R_n$ is orthogonal, left
multiplication by $R_n$ preserves singular values.  Therefore
\[
s_j(C_n(x))=s_j(xI_n-A_n(a)),\qquad j=1,\ldots,n.
\]
Because $C_n(x)$ is real symmetric, its singular values are precisely the
absolute values of its eigenvalues.  The claimed formula follows.
\end{proof}

\section*{Step 2: The inertia calculation}

Put
\[
s=\sqrt a,
\qquad
D_n=\diag(1,s,s^2,\ldots,s^{n-1}),
\]
and
\[
J_n=D_n^{-1}A_n(a)D_n
=
\operatorname{tridiag}(s,0,s).
\]
Thus $J_n$ is the real symmetric Jacobi matrix with zero diagonal and constant
off-diagonal entry $s$.

\begin{lemma}\label{lem:inertia}
Let
\[
I_r=(\mu_{r+1}^{(n)},\mu_r^{(n)}),
\qquad r=1,\ldots,n-1.
\]
For $x\in I_r$,
\[
\nu_-(C_n(x))
=
\left\lfloor\frac n2\right\rfloor
+
\begin{cases}
1,& r\text{ odd},\\
0,& r\text{ even},
\end{cases}
\]
where $\nu_-(M)$ denotes the number of negative eigenvalues of a real symmetric
matrix $M$, counted with multiplicity.  Consequently,
\[
\lambda_{q_n}^{(n)}(x)>0\quad\text{if }r\text{ is even},
\]
and
\[
\lambda_{q_n}^{(n)}(x)<0\quad\text{if }r\text{ is odd}.
\]
Moreover,
\[
\lambda_{q_n}^{(n)}(\mu_j^{(n)})=0,
\qquad j=1,\ldots,n.
\]
\end{lemma}

\begin{proof}
Since
\[
R_nD_n=s^{n-1}D_n^{-1}R_n,
\]
and $A_n(a)=D_nJ_nD_n^{-1}$, we have
\[
C_n(x)
=R_nD_n(xI_n-J_n)D_n^{-1}
=s^{n-1}D_n^{-1}R_n(xI_n-J_n)D_n^{-1}.
\]
Thus $C_n(x)$ is congruent, up to the positive scalar $s^{n-1}$, to
\[
R_n(xI_n-J_n).
\]
Congruence preserves inertia.

The sine vectors
\[
u_k(i)=\sin\frac{ik\pi}{n+1},
\qquad i=1,\ldots,n,
\]
diagonalize $J_n$:
\[
J_nu_k=\mu_k^{(n)}u_k,
\qquad
\mu_k^{(n)}=2s\cos\frac{k\pi}{n+1}.
\]
They are also eigenvectors of $R_n$, because
\[
R_nu_k=(-1)^{k+1}u_k.
\]
Therefore the eigenvalues of $R_n(xI_n-J_n)$ are
\[
(-1)^{k+1}(x-\mu_k^{(n)}),
\qquad k=1,\ldots,n.
\]
If $x\in I_r$, then $x-\mu_k^{(n)}<0$ for $k\leq r$, and
$x-\mu_k^{(n)}>0$ for $k\geq r+1$.  Hence the number of negative signs is
\[
\#\{1\leq k\leq r:k\text{ odd}\}
+
\#\{r+1\leq k\leq n:k\text{ even}\}.
\]
This equals
\[
\left\lceil\frac r2\right\rceil
+
\left(\left\lfloor\frac n2\right\rfloor-
       \left\lfloor\frac r2\right\rfloor\right)
=
\left\lfloor\frac n2\right\rfloor
+
\begin{cases}
1,& r\text{ odd},\\
0,& r\text{ even}.
\end{cases}
\]
The sign statement for $\lambda_{q_n}^{(n)}(x)$ follows immediately from the
definition $q_n=\lfloor n/2\rfloor+1$.

Finally, $A_n(a)$ has the simple eigenvalues $\mu_j^{(n)}$.  Hence
$B_n(\mu_j^{(n)})$ and $C_n(\mu_j^{(n)})$ have nullity one.  The preceding
inertia formula shows that, when $x$ crosses any $\mu_j^{(n)}$, the unique
eigenvalue crossing zero is the ordered eigenvalue with index $q_n$.  Therefore
\[
\lambda_{q_n}^{(n)}(\mu_j^{(n)})=0.
\]
\end{proof}

\section*{Step 3: Simplicity of the smallest singular value in each open gap}

\begin{lemma}[Jacobi inverse formula]\label{lem:jacobi-inverse}
Let
\[
J_n=\operatorname{tridiag}(s,0,s),\qquad s>0,
\]
and let $x=2s\cos\theta$, where $\sin((n+1)\theta)\neq0$.  Then
$xI_n-J_n$ is invertible.  For $i\leq j$,
\[
\bigl((xI_n-J_n)^{-1}\bigr)_{ij}
=
\frac{
 \sin(i\theta)\sin((n+1-j)\theta)
}{
 s\sin\theta\,\sin((n+1)\theta)
},
\]
and the entries for $i\geq j$ are obtained by symmetry.
\end{lemma}

\begin{proof}
Let $\Delta_m$ be the determinant of the leading $m\times m$ principal minor
of $xI_n-J_n$, with $\Delta_0=1$.  The continuant recurrence is
\[
\Delta_m=x\Delta_{m-1}-s^2\Delta_{m-2},
\qquad m\geq2,
\]
with $\Delta_1=x$.  Since $x=2s\cos\theta$, this gives
\[
\Delta_m=s^m\frac{\sin((m+1)\theta)}{\sin\theta}.
\]
Cramer's rule for a symmetric tridiagonal matrix gives, for $i\leq j$,
\[
\bigl((xI_n-J_n)^{-1}\bigr)_{ij}
=s^{j-i}\frac{\Delta_{i-1}\Delta_{n-j}}{\Delta_n}.
\]
Substituting the displayed formula for $\Delta_m$ yields the claim.
\end{proof}

\begin{lemma}\label{lem:green-simplicity}
For every $r=1,\ldots,n-1$ and every
\[
x\in I_r=(\mu_{r+1}^{(n)},\mu_r^{(n)}),
\]
the smallest singular value of $B_n(x)=xI_n-A_n(a)$ is simple:
\[
s_{\min}(B_n(x))<s_2(B_n(x)).
\]
\end{lemma}

\begin{proof}
Fix $x\in I_r$.  Since $x$ is not an eigenvalue of $A_n(a)$, the matrix
$B_n(x)$ is invertible.  By Lemma~\ref{lem:symmetric},
\[
C_n(x)=R_nB_n(x)
\]
is real symmetric.  Also
\[
C_n(x)^{-1}=B_n(x)^{-1}R_n.
\]
Since right multiplication by the orthogonal matrix $R_n$ preserves singular
values,
\[
s_j(B_n(x)^{-1})=s_j(C_n(x)^{-1}),
\qquad j=1,\ldots,n.
\]
Because $C_n(x)^{-1}$ is real symmetric, the singular values of
$C_n(x)^{-1}$ are the absolute values of its eigenvalues.  Thus the largest
singular value of $B_n(x)^{-1}$ is simple if and only if the spectral radius of
$C_n(x)^{-1}$ is attained by a unique simple eigenvalue.

We now identify $C_n(x)^{-1}$ with the Green kernel in
Theorem~\ref{thm:GK}.  Write
\[
x=2s\cos\theta,
\qquad
\theta\in\left(\frac{r\pi}{n+1},\frac{(r+1)\pi}{n+1}\right).
\]
From the identity
\[
C_n(x)
=s^{n-1}D_n^{-1}R_n(xI_n-J_n)D_n^{-1},
\]
we get
\[
C_n(x)^{-1}
=s^{1-n}D_n(xI_n-J_n)^{-1}R_nD_n.
\]
Lemma~\ref{lem:jacobi-inverse} gives, for
$i\leq j$,
\[
\bigl((xI_n-J_n)^{-1}\bigr)_{ij}
=
\frac{
 \sin(i\theta)\sin((n+1-j)\theta)
}{
 s\sin\theta\,\sin((n+1)\theta)
},
\]
and the symmetric formula for $i\geq j$.  Substituting this into the displayed
formula for $C_n(x)^{-1}$ gives exactly the matrix $K_{n,s,\theta}$ from
Theorem~\ref{thm:GK}.  Hence the spectral radius of $C_n(x)^{-1}$ is attained
by a unique simple eigenvalue.

It follows that the largest singular value of $B_n(x)^{-1}$ is simple.  Since
$B_n(x)$ is invertible, the singular values of the inverse are
\[
s_j(B_n(x)^{-1})=\frac{1}{s_{n+1-j}(B_n(x))},
\]
with the singular values of $B_n(x)$ listed increasingly on the right.  Thus
simplicity of the largest singular value of $B_n(x)^{-1}$ is exactly
simplicity of the smallest singular value of $B_n(x)$.  Therefore
\[
s_{\min}(B_n(x))<s_2(B_n(x)).
\]
\end{proof}

\section*{Step 4: The median branch cannot lose the minimum}

\begin{proposition}\label{prop:median-gap}
For every $r=1,\ldots,n-1$ and every $x\in I_r$,
\[
s_{\min}(B_n(x))=
\left|\lambda_{q_n}^{(n)}(x)\right|.
\]
\end{proposition}

\begin{proof}
Fix a gap $I_r=(\mu_{r+1}^{(n)},\mu_r^{(n)})$.  Define
\[
E_r=
\left\{
 x\in I_r:
 s_{\min}(B_n(x))=
 \left|\lambda_{q_n}^{(n)}(x)\right|
\right\}.
\]
We prove that $E_r=I_r$.

First, $E_r$ is nonempty.  Indeed, at each endpoint $\mu_j^{(n)}$ of the gap,
Lemma~\ref{lem:inertia} gives
\[
\lambda_{q_n}^{(n)}(\mu_j^{(n)})=0,
\]
and the nullity of $C_n(\mu_j^{(n)})$ is one.  All other eigenvalues of
$C_n(\mu_j^{(n)})$ are therefore nonzero.  By continuity of the ordered
eigenvalues of a real symmetric matrix, $|\lambda_{q_n}^{(n)}(x)|$ is strictly
smaller than the absolute values of all other eigenvalues for all $x$ in a
sufficiently small one-sided neighbourhood of the endpoint inside $I_r$.  Hence
$E_r\neq\varnothing$.

Second, $E_r$ is closed in $I_r$, because both
$s_{\min}(B_n(x))$ and $|\lambda_{q_n}^{(n)}(x)|$ are continuous functions of
$x$.

Third, $E_r$ is open in $I_r$.  Let $x_0\in E_r$.  By
Lemma~\ref{lem:green-simplicity}, the smallest singular value of $B_n(x_0)$ is
simple.  Since the singular values of $B_n(x_0)$ are the absolute values of the
eigenvalues of $C_n(x_0)$, this means that
$|\lambda_{q_n}^{(n)}(x_0)|$ is strictly smaller than
$|\lambda_j^{(n)}(x_0)|$ for every $j\neq q_n$.  By continuity, the same strict
inequalities persist in a neighbourhood of $x_0$.  Therefore that
neighbourhood is contained in $E_r$, so $E_r$ is open.

The interval $I_r$ is connected.  Since $E_r\subset I_r$ is nonempty, open,
and closed in $I_r$, we have $E_r=I_r$.
\end{proof}

\section*{Conclusion}

\begin{theorem}[Median branch equals the smallest singular value]\label{thm:median}
For every $n\geq2$, every $0<a<1$, and every
\[
x\in[\mu_n^{(n)},\mu_1^{(n)}],
\]
one has
\[
\boxed{
 s_{\min}\bigl(xI_n-A_n(a)\bigr)
 =
 \left|
 \lambda_{\lfloor n/2\rfloor+1}^{(n)}
 \bigl(R_n(xI_n-A_n(a))\bigr)
 \right|.
}
\]
\end{theorem}

\begin{proof}
The assertion holds on each open gap
\[
I_r=(\mu_{r+1}^{(n)},\mu_r^{(n)}),
\qquad r=1,\ldots,n-1,
\]
by Proposition~\ref{prop:median-gap}.  At each endpoint $x=\mu_j^{(n)}$, both
sides are zero: the left-hand side is zero because $\mu_j^{(n)}$ is an
eigenvalue of $A_n(a)$, and the right-hand side is zero by
Lemma~\ref{lem:inertia}.  Hence the equality holds on the whole interval
$[\mu_n^{(n)},\mu_1^{(n)}]$.
\end{proof}

\begin{thebibliography}{9}

\bibitem{GK}
F.~R. Gantmacher and M.~G. Krein,
\emph{Oscillation Matrices and Kernels and Small Vibrations of Mechanical Systems},
AMS Chelsea Publishing, Providence, RI, revised edition, 2002.

\bibitem{Pinkus}
A. Pinkus,
\emph{Totally Positive Matrices},
Cambridge Tracts in Mathematics, Cambridge University Press, 2010.

\end{thebibliography}

\end{document}
